在线策略蒸馏（On-Policy Distillation，OPD）通过学生自生成轨迹上的密集反馈，将教师的推理行为迁移给学生。常见的sampled-token实现为每个生成token分配单独的反馈信号，但一个局部推理决策可能由多个token共同表达。在短片段内共享反馈，能否改善这一学习过程？本文提出Block3 Mean，将连续三个token的师生对数概率差取平均，并将共享信号与联合block重要性比率、按有效长度加权的目标结合。该方法复用已有教师分数，不需要额外奖励模型或过程标注。我们分析其更新动力学，并在不同师生对和保存权重上评测数学推理能力。在历史Qwen3-1.7B对照的Step200，Block3相对sampled-token OPD的Avg@8在MATH500和AMC23上分别提高8.90和5.42个百分点；官方Qwen3-8B到1.7B指令版的Step200也在四项benchmark上获得正向Avg@8差值，而其他配对或保存点可能更有利于逐token更新。我们从局部信用共享、梯度尺度和裁剪几何解释这些比较所涉及的优化变化，并提供区分数值稳定性与生成质量的诊断。这些单seed实验与数学分析刻画了反馈粒度这一推理蒸馏设计选择，并指出值得进一步受控检验的条件。
